\section{Motivation}
\label{sec:motivation}

\subsection{Origin in fluid dynamics}
\label{subsec:fluid-origin}

Burgers' equation was introduced by J. M. Burgers as a simplified model
for turbulence \cite{burgers1948}, isolating the two competing effects
that make the Navier--Stokes equations hard to analyse: nonlinear
advection, which steepens a velocity profile, and viscous diffusion,
which smooths it. The one-dimensional Navier--Stokes momentum equation is
\begin{equation}
  \label{eq:1d-navier-stokes}
  u_t + u\,u_x = -\frac{1}{\rho} p_x + \nu\, u_{xx},
\end{equation}
where $u(x,t)$ is velocity, $\rho$ density, $p$ pressure, and $\nu > 0$
viscosity. In more than one dimension, $p$ is determined implicitly by
the incompressibility constraint $\nabla \cdot u = 0$; this nonlocal
coupling is what makes Navier--Stokes hard to control. In one dimension,
$u_x = 0$ would force $u$ constant, so the constraint is dropped along
with the pressure term entirely, leaving
\begin{equation}
  \label{eq:viscous-burgers}
  u_t + u\, u_x = \nu\, u_{xx},
\end{equation}
the \emph{viscous Burgers' equation}. Setting $\nu = 0$ gives the
\emph{inviscid} equation
\begin{equation}
  \label{eq:inviscid-burgers}
  u_t + u\, u_x = 0,
\end{equation}
which is the main object of study in Sections
\ref{sec:characteristics}--\ref{sec:riemann-problem}, before viscosity is
reintroduced in Section~\ref{sec:extensions}. What remains after this
simplification is a single scalar equation depending only on local values
of $u$, but one that retains the nonlinear term $u\,u_x$ responsible for
the spontaneous formation of discontinuities studied from
Section~\ref{sec:shock-formation} onward.

\begin{remark}[Vorticity-free flow]
\label{rem:vorticity-free}
Without the constraint $\nabla \cdot u = 0$, the flow may be taken
irrotational, so $u = \varphi_x$ for a scalar potential $\varphi$. This is
what allows the Cole--Hopf transformation (Section~\ref{sec:extensions})
to linearise \eqref{eq:viscous-burgers} into the heat equation.
\end{remark}

\subsection{Burgers' equation as a scalar conservation law}
\label{subsec:conservation-law-framework}

Equations \eqref{eq:inviscid-burgers} and \eqref{eq:viscous-burgers} are
instances of a \emph{scalar conservation law},
\begin{equation}
  \label{eq:general-conservation-law}
  u_t + f(u)_x = 0,
\end{equation}
with flux $f$, together with its viscous regularisation $u_t + f(u)_x =
\nu\, u_{xx}$. For Burgers' equation,
\begin{equation}
  \label{eq:burgers-flux}
  f(u) = \frac{u^2}{2},
\end{equation}
so $f(u)_x = u\,u_x$ recovers \eqref{eq:inviscid-burgers}. Such equations
describe a locally conserved quantity transported at state-dependent
speed $f'(u)$, and arise throughout gas dynamics and multiphase flow.
Burgers' equation is the simplest nontrivial case (a linear $f$
gives only the trivial transport equation $u_t + c\,u_x=0$), and is the
standard first example for the general theory of characteristics, shocks,
weak solutions, and entropy conditions developed in
Sections~\ref{sec:characteristics}--\ref{sec:entropy-condition}.

\subsection{Why study Burgers' equation}
\label{subsec:significance}

Three features make Burgers' equation worth studying in its own right.
First, it is the simplest equation whose smooth solutions can lose
smoothness in finite time despite smooth data and no forcing; because it
is scalar and one-dimensional, this can be analysed completely, unlike
the analogous (and in three dimensions still open) questions for Euler
and Navier--Stokes. Second, once weak solutions are admitted, it exhibits
genuine non-uniqueness (Section~\ref{sec:nonuniqueness}), requiring an
entropy condition (Section~\ref{sec:entropy-condition}) to restore
well-posedness --- a phenomenon common to conservation laws in general,
studied here with minimal technical overhead. Third, the viscous equation
\eqref{eq:viscous-burgers} is one of the few nonlinear PDEs solvable in
closed form, via the Cole--Hopf transformation
(Section~\ref{sec:extensions}), allowing a direct, rigorous comparison
between the viscous and inviscid theories.